\documentclass[11pt]{article}

\usepackage{hw_style}

% Set due date for this homework (updated to 2025)
\setduedate{12}{24}

% Setup homework number and header
\hwnumber{10}

\begin{document}

\begin{hwproblem}{}
 请利用大语言模型，根据Sturm-Liouville本征值问题的理论，导出一般的罗德里格斯(Rodrigues)公式:
 $$
 f_n(x)=\frac{a_n}{w(x)} \frac{d^n}{d x^n}\left[w(x)(p(x))^n\right] ,
 $$
 其中$a_n$为与归一化有关的常数. 利用回路积分导出施列夫利(Schlaefli)积分表示:
 $$
 f_n(x)=\frac{1}{w(x)} \frac{n!}{2 \pi i} \oint_C \frac{w(z)(p(z))^n}{(z-x)^{n+1}} d z ,
 $$
 其中$C$为包含点$x$的闭合回路.
\end{hwproblem}


\begin{hwproblem}{}
  厄米多项式(Hermite polynomials) $H_n(x)$的生成函数为
  $$
  e^{-t^2+2 t x}=\sum_{n=0}^{\infty} H_n(x) \frac{t^n}{n !}
  $$
  请利用生成函数推导出下面的递推关系式,
  \begin{itemize}
    \item[(1)] $2 x H_n(x)-2 n H_{n-1}(x)=H_{n+1}(x)$;
    \item[(2)] $2 n H_{n-1}(x)=H_n^{\prime}(x)$ . 
  \end{itemize}
\end{hwproblem}

\vspace{0.5cm}

% \vspace{4em}

\begin{hwproblem}{}
根据勒让德函数的生成函数,推导下面的递推关系式
\begin{itemize}
  \item[(1)] $
  (2 \ell+1) P_\ell(x)=P_{\ell+1}^{\prime}(x)-P_{\ell-1}^{\prime}(x),
  $
  \item[(2)]
  $
  P_{\ell+1}^{\prime}(x)=(\ell+1) P_\ell(x)+x P_\ell^{\prime}(x).
  $
\end{itemize}

\end{hwproblem}

\vspace{0.5cm}
  
% \vspace{4em}

% \begin{hwproblem}{}
%   使用递推关系(recurrence relation)或生成函数或罗德里格斯公式验证连带勒让德函数满足
%   $$
%   \begin{aligned}
%   P_{2 l}^1(0) & =0, \\
%   P_{2 l+1}^1(0) & =(-1)^{l} \frac{(2 l+1) ! !}{(2 l) ! !}.
%   \end{aligned}
%   $$
% \end{hwproblem}

\vspace{0.5cm}




\begin{hwproblem}{}
计算下列积分
\begin{itemize}
  \item[(1)] $$
  \int_{-1}^1\left(1-x^2\right) P_k^{\prime}(x) P_l^{\prime}(x) d x
  $$
  \item[(2)]
  $$
\int_{-1}^1 x P_k(x) P_{k+1}(x) d x
$$
\end{itemize}
\end{hwproblem}

\vspace{0.5cm}

\begin{hwproblem}{}
  求解下列定解问题: 
  $$\left\{\begin{array}{l}
    \nabla^2 u=0, a<r<b, \\ 
    \left.u\right|_{r=a}=u_0,\left.u\right|_{r=b}=u_0 \cos ^2 \theta
  \end{array}\right.$$
\end{hwproblem}

\vspace{0.5cm}


% \begin{hwproblem}{}
%   球谐函数的一般表达式为
%   $$
%   Y_l^m(\theta, \varphi) \equiv \sqrt{\frac{2 l+1}{4 \pi} \frac{(l-m) !}{(l+m) !}} P_l^m(\cos \theta) e^{i m \varphi}
%   $$
%   请求出确定以下的具体表达式$m$的取值范围,写出其中$m\geq 0$关于$\theta,\varphi$的表达式, 并转换成直角坐标系的形式
%   \begin{itemize}
%     \item[(1)] $Y_0^m(\theta,\varphi)$
%     \item[(2)] $Y_1^m(\theta,\varphi)$
%     \item[(3)] $Y_2^m(\theta,\varphi)$
%   \end{itemize}
%   \end{hwproblem}

\vspace{0.5cm}

\hwrequirements

\end{document}